%% kapa: exam=1 round=34 year=2023 period=1 form=A subject=real_estate subject_ko=부동산학원론 seq=5 number=85
%% source: Q-Net 감정평가사 기출문제 (cst00309) artlSeq=5211789 file=raw/1008/5211789/2023년도 제34회 감정평가사 1차시험 1교시.pdf
%% conversion: kordoc 4.24.51 -> tools/gen1.py (auto); review=proofread-sonnet; flags=none
\begin{question}{85}
토지의 특성과 내용에 관한 설명으로 옳지 않은 것은?
\begin{choices}
  \item 토지는 시간의 경과에 의해 마멸되거나 소멸되지 않으므로 투자재로서 선호도가 높다.
  \item 물리적으로 완전히 동일한 토지는 없으므로 부동산시장은 불완전경쟁시장이 된다.
  \item 토지는 공간적으로 연결되어 있으므로 외부효과를 발생시키고, 개발이익 환수의 근거가 된다.
  \item 토지는 용익물권의 목적물로 활용할 수 있으므로 하나의 토지에 다양한 물권자가 존재할 수 있다.
  \item 토지의 소유권은 정당한 이익있는 범위내에서 토지의 상하에 미치며, 한계고도와 한계심도의 범위는 법률로 정하고 있다.
\end{choices}
\end{question}
\answer{5}
